\section{Problem Analysis}
For transcription of voice commands, a reliable and accurate speech recognition system is required. For this an AI solution is prefered due to the complexity and variability of human speech. OpenAI's Whisper model is chosen for this task, as it has the required capabilities and is open source, allowing for easy integration and customization. It has also scored as the high ends of speech recognition benchmarks, between different models testing their WER (Word Error Rate).

For sound localisation, a method to determine the direction of the sound source is needed. This requires multiple microphones set apart from one another. By using 3 microphones and comparing signals and finding the phase difference between them, the direction of the sound source can be estimated. For this the GCC-PHAT (Generalized Cross-Correlation with Phase Transform) algorithm is chosen, as it will offer robust performance in noisy environments. Using the speed of sound the position can then be estimated by triangulating from the phase differences.

As sound localisation needs multiple microphones, with the same clock to ensure synchronized sampling I2S microphones are used as they share a clock. The Raspberry Pi cannot handle more than one stereo channel of I2S (Inter-Integrated Circuit Sound) audio, so an ESP32 microcontroller will be needed to interface with the microphones and send the audio data as it has 2 stereo I2S interfaces, totalling 4 mono channels. The microphones can then be mounted on a 3D printed structure which can be used for triangulation.

Whilst Raspberry Pi and ESP32 are both on the robot, neither have the computational power to run the Whisper model in real-time. Therefore a seperate PC will be needed to handle the transcription and sound localisation processing.


Communication between the PC and the Raspberry Pi will be handled using MQTT (Message Queuing Telemetry Transport), a lightweight Pub-Sub messaging protocol as it was already in use for command transfer to the robot. The ESP32 do have a Wifi module and can therefore also use MQTT to send the audio date, but the efficiency and power of the wifi module on the ESP32 might not be enough to handle the audio streaming over wifi. If that is the case the both the ESP32 and the Raspberry pi is also mounted with different wired communication protocols allowing to transfer the audio data, where the Raspberry Pi then can send it over MQTT to the PC. While dedicated audio streaming protocols like RTP, RTSP or WebRTC exist, MQTT's simplicity and low overhead make it suitable for this project's needs.